% Options for packages loaded elsewhere
\PassOptionsToPackage{unicode}{hyperref}
\PassOptionsToPackage{hyphens}{url}
\PassOptionsToPackage{dvipsnames,svgnames,x11names}{xcolor}
\documentclass[
  11pt,
  letterpaper,
]{article}
\usepackage{xcolor}
\usepackage[margin=0.88in]{geometry}
\usepackage{amsmath,amssymb}
\setcounter{secnumdepth}{-\maxdimen} % remove section numbering
\usepackage{iftex}
\ifPDFTeX
  \usepackage[T1]{fontenc}
  \usepackage[utf8]{inputenc}
  \usepackage{textcomp} % provide euro and other symbols
\else % if luatex or xetex
  \usepackage{unicode-math} % this also loads fontspec
  \defaultfontfeatures{Scale=MatchLowercase}
  \defaultfontfeatures[\rmfamily]{Ligatures=TeX,Scale=1}
\fi
\usepackage{lmodern}
\ifPDFTeX\else
  % xetex/luatex font selection
\fi
\usepackage{titlesec}
\usepackage{fancyhdr}
\definecolor{primaryblue}{RGB}{2,100,160}
\definecolor{accentindigo}{RGB}{60,50,190}
\definecolor{darkslate}{RGB}{28,38,52}
\definecolor{boxborder}{RGB}{210,220,230}
\titleformat{\section}{\Large\bfseries\color{primaryblue}}{}{0pt}{}[\vspace{-0.3em}{\color{boxborder}\hrule height 0.7pt}]
\titleformat{\subsection}{\large\bfseries\color{accentindigo}}{}{0pt}{}
\titleformat{\subsubsection}{\normalsize\bfseries\color{darkslate}}{}{0pt}{}
\titlespacing*{\section}{0pt}{5pt plus 1.5pt minus 1pt}{2.5pt plus 1pt minus 1pt}
\titlespacing*{\subsection}{0pt}{3.5pt plus 1pt minus 1pt}{1.5pt plus 1pt minus 1pt}
\pagestyle{fancy}
\fancyhf{}
\fancyhead[L]{\small\sffamily\color{darkslate}\textbf{Cultural Mechanics}}
\fancyhead[R]{\small\sffamily\color{primaryblue}\textbf{Lecture 06}}
\fancyfoot[L]{\scriptsize\sffamily\color{darkslate}Semantic Mechanics Project $\cdot$ Mathematical Linguistics Group}
\fancyfoot[R]{\small\sffamily\color{darkslate}\thepage}
\renewcommand{\headrulewidth}{0.4pt}
\renewcommand{\footrulewidth}{0.4pt}
\setlength{\headheight}{14pt}
% Use upquote if available, for straight quotes in verbatim environments
\IfFileExists{upquote.sty}{\usepackage{upquote}}{}
\IfFileExists{microtype.sty}{% use microtype if available
  \usepackage[]{microtype}
  \UseMicrotypeSet[protrusion]{basicmath} % disable protrusion for tt fonts
}{}
\makeatletter
\@ifundefined{KOMAClassName}{% if non-KOMA class
  \IfFileExists{parskip.sty}{%
    \usepackage{parskip}
  }{% else
    \setlength{\parindent}{0pt}
    \setlength{\parskip}{6pt plus 2pt minus 1pt}}
}{% if KOMA class
  \KOMAoptions{parskip=half}}
\makeatother
\usepackage{longtable,booktabs,array}
\newcounter{none} % for unnumbered tables
\usepackage{calc} % for calculating minipage widths
% Correct order of tables after \paragraph or \subparagraph
\usepackage{etoolbox}
\makeatletter
\patchcmd\longtable{\par}{\if@noskipsec\mbox{}\fi\par}{}{}
\makeatother
% Allow footnotes in longtable head/foot
\IfFileExists{footnotehyper.sty}{\usepackage{footnotehyper}}{\usepackage{footnote}}
\makesavenoteenv{longtable}
\setlength{\emergencystretch}{3em} % prevent overfull lines
\providecommand{\tightlist}{%
  \setlength{\itemsep}{0pt}\setlength{\parskip}{0pt}}
\usepackage{bookmark}
\IfFileExists{xurl.sty}{\usepackage{xurl}}{} % add URL line breaks if available
\urlstyle{same}
\hypersetup{
  colorlinks=true,
  linkcolor={blue},
  filecolor={Maroon},
  citecolor={Blue},
  urlcolor={teal},
  pdftitle={Lecture 06: Greed, Toxicity, and Repair},
  pdfauthor={Semantic Mechanics Project},
  pdfcreator={LaTeX via pandoc}}

\author{}
\date{}

\begin{document}

\section{Cultural Mechanics: Lecture
06}\label{cultural-mechanics-lecture-06}

\textbf{SEMANTIC MECHANICS PROJECT · MATHEMATICAL LINGUISTICS GROUP}

\begin{center}\rule{0.5\linewidth}{0.5pt}\end{center}

\subsection{Lecture 06: Greed, Toxicity, and
Repair}\label{lecture-06-greed-toxicity-and-repair}

\subsubsection{Listening, Grief, Resource Defense, and the Political
Economy of
Harm}\label{listening-grief-resource-defense-and-the-political-economy-of-harm}

\begin{itemize}
\tightlist
\item
  \textbf{Course:} Cultural Mechanics
\item
  \textbf{Initiative:} The Semantic Mechanics Project
\item
  \textbf{Term:} Autumn 2026
\item
  \textbf{Prerequisites:} Social Psychology, Behavioral Economics,
  Cultural History, and Introductory Epidemiology
\item
  \textbf{Compiled lecture PDF:}
  \href{lecture_06_greed_toxicity_and_repair.pdf}{\texttt{lecture\_06\_greed\_toxicity\_and\_repair.pdf}}
\item
  \textbf{LaTeX source:}
  \href{lecture_06_greed_toxicity_and_repair.tex}{\texttt{lecture\_06\_greed\_toxicity\_and\_repair.tex}}
\item
  \textbf{Content note:} This lecture discusses sexual assault, family
  estrangement, grief, racism, enslavement, and economic coercion. No
  student is expected to disclose personal experience.
\item
  \textbf{Foundational readings:}

  \begin{itemize}
  \tightlist
  \item
    DeKornfeld, Aviva (host). (2022). ``902: The Missing Minutes.''
    \emph{This American Life}, Act One: ``Known Unknowns.''
    \href{https://www.thisamericanlife.org/902/transcript}{Transcript}.
  \item
    Bureau for Listening. ``Aggressive Listening.''
    \href{https://bureauforlistening.com/agressive-listening/}{Project
    description}.
  \item
    \emph{Unlearning Insecurity: Active Inference, Bayesian Precision
    Weighting, and Cognitive Restructuring.}
    \href{../../../hpc/chaoshearing/notebooks/unlearning_insecurity_active_inference.ipynb}{Companion
    notebook}.
  \item
    Kupers, Terry A. (2005). ``Toxic masculinity as a barrier to mental
    health treatment in prison.'' \emph{Journal of Clinical Psychology},
    61(6), 713--724.
  \item
    Connell, R. W. (1995). \emph{Masculinities}. University of
    California Press.
  \item
    Kahneman, Daniel, and Amos Tversky. (1979). ``Prospect Theory: An
    Analysis of Decision under Risk.'' \emph{Econometrica}, 47(2),
    263--291.
  \item
    Rothman, Kenneth J., Sander Greenland, and Timothy L. Lash. (2008).
    \emph{Modern Epidemiology}, 3rd ed.~Lippincott Williams \& Wilkins.
  \end{itemize}
\end{itemize}

\begin{center}\rule{0.5\linewidth}{0.5pt}\end{center}

\subsection{1. A Question About What We
Keep}\label{a-question-about-what-we-keep}

Greed is often presented as a private defect: one person wants too much.
But what counts as ``too much,'' who gets to define enough, and what
happens to everyone else when one person or institution accumulates
control? This lecture follows a change in scale:

\begin{enumerate}
\def\labelenumi{\arabic{enumi}.}
\tightlist
\item
  A story of assault and the struggle to know what happened.
\item
  Listening as attention, pressure, and a contest over who controls the
  subject.
\item
  Grief as a relation to loss rather than a failure to recover.
\item
  Resource defense in animals, described without projecting human moral
  categories onto them.
\item
  Epidemiology and behavioral economics as tools for mapping and
  preventing harm.
\item
  Class, racial hierarchy, debt, reparations, reconstruction, and basic
  income as questions of institutional design.
\end{enumerate}

The through-line is \textbf{resource security}: bodies, attention,
credibility, money, land, time, and the ability to refuse. Greed is one
possible name for harmful accumulation, but the analysis must ask what
is accumulated, by whom, through what rules, and at whose expense.

\subsubsection{Learning objectives}\label{learning-objectives}

By the end of the lecture, students should be able to:

\begin{itemize}
\tightlist
\item
  Explain how ``toxic masculinity'' names harmful norms and reward
  systems, not an essence shared by men.
\item
  Analyze listening as a practice with power, boundaries, and
  accountability.
\item
  Distinguish ordinary grief from a clinical disorder and from social
  abandonment.
\item
  Analyze resource defense and cooperation as behaviors humans share
  with other animals, while accounting for the distinctive reach of
  human institutions and culture.
\item
  Distinguish fixed-sum competition from positive-sum cooperation, and
  ask who benefits when the total surplus grows.
\item
  Explain why a convincing cloned voice is not proof of identity,
  consent, or viewpoint, and apply practical verification steps.
\item
  Use an epidemiologic frame to study distributions and causes of harm
  without calling greed, racism, or grief a ``disease of the mind.''
\item
  Apply behavioral-economics tools to prevention while identifying when
  ``nudging'' becomes paternalism or surveillance.
\item
  Evaluate debt relief, reparations, reconstruction, and universal basic
  income as distinct approaches to repair and security.
\end{itemize}

\begin{center}\rule{0.5\linewidth}{0.5pt}\end{center}

\subsection{2. ``Toxic'': From Poison to a Social
Diagnosis}\label{toxic-from-poison-to-a-social-diagnosis}

The word \emph{toxic} evokes a substance that harms through exposure. In
contemporary social use, it often describes a relationship, norm,
institution, or environment that causes damage. The metaphor is vivid
but can obscure causation: substances have doses and pathways; social
harms have histories, incentives, power relations, and people who can
choose differently.

``Toxic masculinity'' has not had one stable meaning. The phrase
circulated in late-twentieth-century men's-movement discourse, including
mythopoetic writing, and later became common in feminist, psychological,
and popular debate. In present research and practice, it is most useful
when it describes \textbf{restrictive norms}---for example, entitlement,
dominance, emotional suppression, sexual coercion, or the treatment of
vulnerability as weakness---and the social rewards that can sustain
them. It is not a synonym for masculinity, men, or male identity.
Masculinities vary across history and culture; people can resist,
revise, or reproduce norms.

The distinction matters in the 2022 \emph{This American Life} episode,
Act One, ``Known Unknowns,'' which tells the story of Vienna Luke
seeking an account of an assault by her brother-in-law. The episode
describes her memory gaps, the family's prolonged avoidance, and her
eventual, police-supported recorded meeting in which Justin acknowledges
drugging her, removing her clothes, photographing her, and other acts.
The story is not offered as a template for survivors: the meeting was an
unusually supported investigative operation, and its result does not
imply that a survivor must confront an assailant, obtain a confession,
report to police, or pursue closure.

Justin's descriptions of ``lines'' he believed he had not crossed and
his detours into the general problem of pornography illustrate a crucial
distinction between \textbf{self-image} and \textbf{accountability}. A
person can preserve a story of being good by defining harm narrowly,
treating intent as exculpatory, or replacing a specific act with a broad
cultural explanation. Explanation can illuminate conditions; it does not
erase agency or responsibility.

The phrase ``toxic masculinity'' can help name a pattern, but it can
also become a shortcut that stigmatizes men as a group, overlooks
differences in power, or mistakes an explanation for a solution. A
careful analysis identifies:

{\def\LTcaptype{none} % do not increment counter
\begin{longtable}[]{@{}
  >{\raggedright\arraybackslash}p{(\linewidth - 2\tabcolsep) * \real{0.5000}}
  >{\raggedright\arraybackslash}p{(\linewidth - 2\tabcolsep) * \real{0.5000}}@{}}
\toprule\noalign{}
\begin{minipage}[b]{\linewidth}\raggedright
Question
\end{minipage} & \begin{minipage}[b]{\linewidth}\raggedright
What to examine
\end{minipage} \\
\midrule\noalign{}
\endhead
\bottomrule\noalign{}
\endlastfoot
\textbf{Norm} & What behavior is framed as masculine, respectable, or
entitled? \\
\textbf{Reward} & Who gains approval, status, access, money, or impunity
for following it? \\
\textbf{Cost} & Who bears bodily, emotional, social, or economic
harm? \\
\textbf{Institution} & What procedure or power imbalance protects the
behavior? \\
\textbf{Alternative} & What norm, relationship, or accountability
practice could change the outcome? \\
\end{longtable}
}

This keeps the object of critique precise: a harmful practice and its
enabling conditions, not an entire identity category.

\begin{center}\rule{0.5\linewidth}{0.5pt}\end{center}

\subsection{3. Aggressive Listening: Attention, Territory, and
Consent}\label{aggressive-listening-attention-territory-and-consent}

The Bureau for Listening's project frames ``aggressive listening'' as a
kind of \textbf{auditory staring}: sustained attention that can affect,
challenge, or contest the territory of what may be said. It deliberately
resists a closed definition. The useful tension is that listening is
often described as passive or benevolent, while attention can also be
intrusive, demanding, and politically consequential.

In the episode, Vienna repeatedly returns the conversation to the
question of what happened when Justin shifts toward abstractions. Her
insistence is not a general listening technique to prescribe to
survivors; it is an action made possible by preparation, a clear goal,
and a police team monitoring the operation. Her listening is active
because she controls the question and refuses the other speaker's
attempts to change the frame.

We can describe the exchange as a \textbf{frame-control problem}:

\[
\text{Accountability} =
\text{specific act}
+ \text{specific agent}
+ \text{specific consequence}
+ \text{response to the harmed person}
\]

When any term is replaced by a vague generality---``culture,''
``addiction,'' ``mistakes,'' ``everyone does it''---the account may lose
the actor and the consequence. Listening closely means hearing both what
is said and what a speaker's framing leaves out.

But aggressive listening has an ethical edge. Attention can be care; it
can also be extraction. A listener should ask:

\begin{itemize}
\tightlist
\item
  Who set the terms of this conversation?
\item
  Is the speaker safe to stop, pause, or decline?
\item
  Is the listener seeking understanding, evidence, confession, or
  spectacle?
\item
  Who controls the recording, retelling, and circulation of the account?
\item
  What responsibilities follow from hearing?
\end{itemize}

The episode's hidden recording is part of a specific police-supported
investigation. It should not be generalized into advice to secretly
record people. Recording law, safety risks, consent, and an individual's
goals vary. A listener's obligation is not to force disclosure; it is to
respect agency, avoid replacing the speaker's account with speculation,
and respond to concrete safety needs.

\subsubsection{Voice doppelgängers: listening is not
authentication}\label{voice-doppelguxe4ngers-listening-is-not-authentication}

A \textbf{voice doppelgänger} is a generated or edited voice that
imitates a person's vocal identity. It may make words sound as if they
came from someone who never said them, or make a genuine recording
appear to express a perspective the speaker did not hold. This matters
for aggressive listening: sustained attention to an audio artifact is
not the same as listening to the person whose voice it imitates. A voice
can prompt intimacy, fear, or trust, but timbral resemblance alone
cannot establish identity, authorship, consent, context, or intent.

Authorized voice synthesis can also support augmentative communication,
preserve a person's chosen voice when speech is changing, enable
translation, and serve creative performance. Those benefits make
meaningful consent and continuing control important, not optional.

The ethical questions concern more than whether a model can imitate a
voice:

{\def\LTcaptype{none} % do not increment counter
\begin{longtable}[]{@{}
  >{\raggedright\arraybackslash}p{(\linewidth - 2\tabcolsep) * \real{0.5000}}
  >{\raggedright\arraybackslash}p{(\linewidth - 2\tabcolsep) * \real{0.5000}}@{}}
\toprule\noalign{}
\begin{minipage}[b]{\linewidth}\raggedright
Issue
\end{minipage} & \begin{minipage}[b]{\linewidth}\raggedright
Questions for creators, platforms, and listeners
\end{minipage} \\
\midrule\noalign{}
\endhead
\bottomrule\noalign{}
\endlastfoot
\textbf{Permission} & Did the person knowingly agree to the use of their
voice, for this purpose and audience? \\
\textbf{Scope} & Does permission cover only a specific, reviewed
recording, or open-ended generation of new statements? \\
\textbf{Disclosure} & Can listeners tell that speech is synthetic or
materially edited? \\
\textbf{Control} & Can the person review, limit, or revoke future use?
What happens to stored recordings and models? \\
\textbf{Power and harm} & Who is paid, impersonated, made vulnerable, or
burdened with disproving the output? \\
\textbf{After death} & Who may authorize a replica, and how are the
person's wishes and the interests of family and audience represented? \\
\end{longtable}
}

Consent to contribute a voice sample is not consent to every sentence
that can be generated from it. Nor does an authentic recording guarantee
an authentic perspective: it may be edited, taken out of context, or
spoken under pressure. For high-stakes listening, ask who made the
recording, how it reached you, whether it is labeled, and what
independent evidence supports the claim. Treat a purported voice message
as an artifact to verify---not as direct testimony---until its source
and context are established.

\textbf{Practical protection when a voice message requests money,
access, or secrecy:}

\begin{enumerate}
\def\labelenumi{\arabic{enumi}.}
\tightlist
\item
  Pause. Urgency, secrecy, and instructions not to check with anyone are
  warning signs; do not transfer money or share passwords, identity
  documents, or one-time security codes based on a voice alone.
\item
  Verify through a channel you already trust: call a number saved
  independently or contact the person through a separate, established
  account. Do not rely on the incoming number, caller ID, or a link
  supplied in the message.
\item
  If appropriate, use a family check-in phrase agreed in advance, but
  treat it as one signal---not an infallible test. Do not use a
  memorable phrase as an account password or authentication code.
\item
  For financial or account changes, use the institution's official app
  or contact information and ask for time to verify. If money or
  credentials were already sent, contact the bank or service promptly
  through its official channel.
\item
  Preserve the original message, sender details, time, and surrounding
  context; report suspected impersonation to the relevant platform and,
  where appropriate, local authorities. Avoid forwarding the clip as if
  it were authentic.
\end{enumerate}

Voice samples already online cannot be made un-cloneable, and there is
no reliable ``listen harder'' test that proves a voice is genuine.
Automated detectors and audible artifacts can help triage but may miss
sophisticated edits or flag real speech. Provenance labels and
watermarks can add context when present, but their absence does not
prove authenticity and their presence does not prove that the content is
truthful. The core defense is layered verification and a norm that no
one should be pressured to treat a voiceprint as consent or proof.

\begin{center}\rule{0.5\linewidth}{0.5pt}\end{center}

\subsection{4. Insecurity, Active Inference, and the Limits of the
Model}\label{insecurity-active-inference-and-the-limits-of-the-model}

The companion notebook models insecurity as a strong negative prior that
dominates ambiguous evidence. It sketches how avoidance can preserve a
belief by reducing opportunities to gather new information, and how
supported behavioral experiments may provide disconfirming evidence.
This is a useful computational analogy for some patterns of anxious
prediction:

\[
\text{posterior belief} \propto
\text{prior belief} \times \text{likelihood of observed evidence}
\]

An agent who expects rejection may interpret ambiguous social cues as
confirmation, then avoid the situations that could test that
expectation. The notebook's language of precision weighting and active
sampling offers a way to ask how beliefs are maintained and updated.

However, a model is not a diagnosis, and several distinctions are
essential:

\begin{enumerate}
\def\labelenumi{\arabic{enumi}.}
\tightlist
\item
  \textbf{A survivor's uncertainty is not necessarily an inaccurate
  belief.} Memory gaps after trauma are not proof that an assault did or
  did not happen. The responsibility for an assault lies with the person
  who committed it, not with the survivor's inference process.
\item
  \textbf{Evidence is not always safe to seek.} In an abuse or coercive
  relationship, approaching the source of uncertainty can increase
  danger. ``Exposure'' is not a universal prescription.
\item
  \textbf{A social environment can withhold information or credibility.}
  Not every locked belief is maintained by an individual's avoidance;
  institutions and other people can actively prevent access to evidence
  or recognition.
\item
  \textbf{The notebook's equations are explanatory simplifications.} Its
  Kalman-gain and free-energy diagrams do not by themselves validate a
  clinical treatment, establish a universal brain mechanism, or capture
  the social history of trauma.
\end{enumerate}

The best transfer is modest: uncertainty can be worsened by avoidance,
but not all uncertainty is self-generated; gathering information can
support learning, but only when it is safe, wanted, and proportionate.
In Vienna's case, preparation, law enforcement, evidence, and her own
deliberate questions mattered. The episode also shows that facts did not
automatically restore the relationship she hoped to repair.

\begin{center}\rule{0.5\linewidth}{0.5pt}\end{center}

\subsection{5. Grief: Loss Without a
Deadline}\label{grief-loss-without-a-deadline}

Grief is a response to loss, not an error that must be corrected. It may
follow death, but also the loss of trust, bodily safety, a relationship,
a future one expected, or a family's former shape. The episode holds
several kinds of grief at once: Vienna's grief over what was done to
her, the family's altered rituals, and the hoped-for return of closeness
with Sunny.

Grief does not proceed in a universal sequence. Stage models can be
evocative descriptions, but they are not a timetable that every grieving
person must complete. People may move among anger, numbness, yearning,
relief, and ordinary daily life; they may continue a bond with the
person or life they lost. A continuing bond is not necessarily denial,
and a public or legal outcome does not guarantee emotional repair.

Clinical diagnosis has a narrower role. Prolonged grief disorder
describes a persistent, impairing grief response that meets specific
criteria in a particular diagnostic system. It should not be used to
pathologize ordinary sadness, cultural mourning, or an inconvenient
demand for accountability. Nor should a lack of visible grief be treated
as evidence that no harm occurred.

For discussion, separate three questions:

\begin{itemize}
\tightlist
\item
  \textbf{What was lost?} Name the relationship, safety, time, income,
  community, or future.
\item
  \textbf{What support is available?} Consider chosen family, material
  resources, ritual, therapy, legal advocacy, and time.
\item
  \textbf{What remains unresolved because another person or institution
  will not act?} Grief may coexist with a claim for justice; it is not a
  substitute for it.
\end{itemize}

``Closure'' is often imagined as an internal finish line. A more humane
goal may be increased safety, room to remember or not remember, the
ability to set boundaries, and a life not organized entirely around the
harm.

\begin{center}\rule{0.5\linewidth}{0.5pt}\end{center}

\subsection{6. Resource Defense, Cooperation, and Zero-Sum Thinking
Across Animal
Life}\label{resource-defense-cooperation-and-zero-sum-thinking-across-animal-life}

Humans are animals. We share with other species basic needs and evolved
capacities for competition, cooperation, attachment, defense, and
exchange. Human beings also defend food, shelter, territory, offspring,
and valued social positions. Those behaviors are shaped---not simply
dictated---by ecology, development, culture, and institutions. Other
animals likewise cooperate, share, form alliances, and negotiate access.
Neither competition nor cooperation belongs exclusively to humans.

In behavioral ecology, researchers might ask what resource is at stake,
how contest costs compare to likely benefits, whether a resource is
defendable, and what signals or escalation strategies animals use. A
small lizard defending a basking site and a human landlord enforcing a
lease may both involve access and exclusion, but the settings,
capacities, and meanings differ. ``Greed'' is a moral and cultural
interpretation, not a neutral zoological diagnosis.

The comparison is useful when it helps us ask about resource access,
conflict costs, and cooperation without making a natural fact out of a
political arrangement:

{\def\LTcaptype{none} % do not increment counter
\begin{longtable}[]{@{}
  >{\raggedright\arraybackslash}p{(\linewidth - 4\tabcolsep) * \real{0.3333}}
  >{\raggedright\arraybackslash}p{(\linewidth - 4\tabcolsep) * \real{0.3333}}
  >{\raggedright\arraybackslash}p{(\linewidth - 4\tabcolsep) * \real{0.3333}}@{}}
\toprule\noalign{}
\begin{minipage}[b]{\linewidth}\raggedright
Animal-behavior question
\end{minipage} & \begin{minipage}[b]{\linewidth}\raggedright
Human institutional analogue
\end{minipage} & \begin{minipage}[b]{\linewidth}\raggedright
Important difference
\end{minipage} \\
\midrule\noalign{}
\endhead
\bottomrule\noalign{}
\endlastfoot
Is a resource scarce, clumped, or renewable? & How are housing, food,
land, or income distributed? & Scarcity may be socially produced or
politically maintained. \\
What is the energetic or injury cost of defense? & What bodily,
emotional, or civic costs accompany competition? & Human rules can
redistribute costs and authorize coercion. \\
Who can access the resource or signal ownership? & Who has legal title,
bargaining power, or recognized claims? & Human ownership is constituted
by institutions and history. \\
What forms of cooperation or sharing are stable? & Can people create
reciprocal arrangements, public goods, or shared infrastructure? & Human
cooperation can be organized across large populations and
generations. \\
\end{longtable}
}

\subsubsection{Fixed-sum contests and ``why can't we grow the
pie?''}\label{fixed-sum-contests-and-why-cant-we-grow-the-pie}

A \textbf{zero-sum game}, in the strict sense, has a fixed total payoff:
one participant's gain is another's loss. Some contests are genuinely
fixed-sum over a particular scarce resource---for example, one
indivisible object can be awarded to only one claimant. Many social and
economic situations are not fixed-sum. Cooperation, technology, public
health, and shared infrastructure can create \textbf{positive-sum
gains}, where total benefits grow.

But ``grow the pie'' is not a complete answer to distribution. Total
output can expand while ownership and bargaining power remain
concentrated; the gains may flow upward while pollution, unpaid care,
animal suffering, or precarious labor are treated as external costs.
Conversely, redistributing a fixed resource does not necessarily create
new resources, but may be needed for fairness and security. Ask both:

\begin{itemize}
\tightlist
\item
  \textbf{Can the total surplus grow?} What cooperation, innovation, or
  public investment could make more value or reduce avoidable loss?
\item
  \textbf{How is the surplus distributed?} Who owns the assets, captures
  the gains, bears the risks, and has the power to refuse?
\item
  \textbf{What is being counted as a cost?} Does the ``larger pie''
  depend on unpriced environmental damage, coerced labor, or harm to
  other animals?
\end{itemize}

Human institutions extend resource defense beyond immediate physical
contests. Property law, inheritance, credit, labor markets, and
political voice can store claims across time and make exclusion durable.
Cumulative culture can also support reciprocal exchange, collective
care, and rules that limit domination. These capacities do not make
humans categorically separate from nature; they make our social
arrangements unusually consequential and open to redesign.

An animal analogy should therefore illuminate, not excuse. ``Competition
made me do it'' naturalizes a choice; ``people are just territorial''
naturalizes an institution. The better questions are what conditions
elicit defensive behavior, when cooperation is possible, and who can
change the rules.

\begin{center}\rule{0.5\linewidth}{0.5pt}\end{center}

\subsection{7. Epidemiology Without Calling Society
Sick}\label{epidemiology-without-calling-society-sick}

Epidemiology studies the distribution and determinants of health-related
states in populations. Its methods can help investigate the
\textbf{patterning of harm}: where exposure occurs, who is more likely
to experience it, which institutions amplify risk, and what
interventions reduce incidence or severity. They do not establish that
greed, racism, masculinity, or grief are mental diseases.

A population-health frame might map:

\[
\text{conditions and exposures}
\rightarrow \text{patterns of harm}
\rightarrow \text{unequal outcomes}
\rightarrow \text{prevention and repair}
\]

For example, researchers could study how housing insecurity, workplace
power, social isolation, access to care, or institutional impunity
correlate with violence and distress. These are associations to
investigate, not proof that any single exposure causes an individual
act. Avoid confusing a risk factor with a diagnosis or blaming people
who experience the outcome.

The World Health Organization describes social determinants as the
conditions in which people are born, grow, live, work, and age,
including access to power, money, and resources. That framing shifts
inquiry from ``What is wrong with these people?'' to ``Which conditions
distribute risk and protection, and how can they change?'' The metaphor
of a ``social pathogen'' can sometimes draw attention to contagious
norms or repeated institutional practices, but it can also medicalize
dissent, stigmatize groups, and hide agents who benefit. Prefer direct
language---violence, coercion, exclusion, exploitation, or unequal
exposure---when it names the mechanism more accurately.

An ethical population analysis should state its outcome, population,
timeframe, comparison group, and limits; distinguish correlation from
cause; avoid treating group averages as individual destinies; and
involve affected communities in interpreting results. The aim is
prevention and equitable protection, not surveillance or labeling a
community as diseased.

\begin{center}\rule{0.5\linewidth}{0.5pt}\end{center}

\subsection{8. Behavioral Economics: Steering Away from
Harm}\label{behavioral-economics-steering-away-from-harm}

Behavioral economics asks how decisions are shaped by attention,
defaults, framing, habits, limited time, and uncertainty---not just
stable preferences and perfect information. Prospect theory, for
example, describes how choices under risk can depend on reference points
and losses. Scarcity can narrow attention toward urgent demands, while
unequal bargaining power can make a nominal ``choice'' costly to refuse.

These insights can inform harm prevention:

\begin{itemize}
\tightlist
\item
  \textbf{Change defaults:} Make consent, privacy, and safe reporting
  the default rather than an extra burden.
\item
  \textbf{Reduce friction:} Offer accessible routes to housing, health
  care, legal advice, and income support.
\item
  \textbf{Make consequences legible:} State who is accountable, what
  will happen next, and where a complaint goes.
\item
  \textbf{Design for repair:} Provide predictable, timely restitution
  and accessible non-carceral supports.
\item
  \textbf{Protect meaningful choice:} Preserve opt-out, confidentiality,
  due process, and avenues to contest institutional decisions.
\end{itemize}

Yet a ``nudge'' is not automatically benevolent. A default can help
people claim a benefit, or quietly enroll them in surveillance. A
reminder can improve safety, or shift responsibility from a powerful
institution to an individual with less power. Behavioral tools should be
judged by who designed them, who benefits, who bears error, whether
people can refuse, and whether the intervention changes the harmful
environment or merely teaches vulnerable people to navigate it.

\subsubsection{The cost of harm}\label{the-cost-of-harm}

Accounting for harm requires more than a price tag. A cost map can
include:

{\def\LTcaptype{none} % do not increment counter
\begin{longtable}[]{@{}
  >{\raggedright\arraybackslash}p{(\linewidth - 2\tabcolsep) * \real{0.5000}}
  >{\raggedright\arraybackslash}p{(\linewidth - 2\tabcolsep) * \real{0.5000}}@{}}
\toprule\noalign{}
\begin{minipage}[b]{\linewidth}\raggedright
Harm domain
\end{minipage} & \begin{minipage}[b]{\linewidth}\raggedright
Examples to consider
\end{minipage} \\
\midrule\noalign{}
\endhead
\bottomrule\noalign{}
\endlastfoot
\textbf{Physical} & Injury, pain, disability, medical care, lost
mobility \\
\textbf{Cognitive and emotional} & Fear, sleep disruption, impaired
concentration, grief, loss of trust \\
\textbf{Social} & Isolation, family rupture, lost education, reduced
civic participation \\
\textbf{Economic} & Unpaid care, missed work, lower wages, legal costs,
displacement \\
\textbf{Ecological and interspecies} & Habitat loss, animal injury,
disrupted food systems \\
\textbf{Institutional} & Investigation, court processes, reduced trust,
resources diverted from prevention \\
\end{longtable}
}

These domains interact. A person may face a physical injury, cognitive
load, lost wages, and isolation at once. Monetizing some costs can help
compare policy choices, but it cannot fully represent dignity, consent,
or irreplaceable loss. A responsible analysis shows what its accounting
leaves out and never makes a harmed person responsible for proving that
their pain is expensive enough to matter.

\begin{center}\rule{0.5\linewidth}{0.5pt}\end{center}

\subsection{9. Class, Race, and the Unequal Distribution of
Credibility}\label{class-race-and-the-unequal-distribution-of-credibility}

All humans belong to one species, \emph{Homo sapiens}. ``Race'' is not a
set of discrete biological subspecies, but racial categories have been
historically made and politically enforced; they have real consequences
because institutions distribute safety, wealth, credibility, and
opportunity through them. Saying ``we are all one human race'' can
affirm shared humanity, but it cannot by itself explain or remedy
racialized inequality.

Class describes relations to resources, work, ownership, and bargaining
power; it is not interchangeable with race. The categories intersect and
vary across place and history. A class analysis that ignores racialized
dispossession can erase how wealth was accumulated. A racial analysis
that ignores labor and property can miss how class power reproduces
inequality. Neither should assume that every person in a group has the
same experience.

The episode's family conflict makes credibility visible at an intimate
scale: who is believed, whose comfort is prioritized, and who is
expected to keep the family together. At a societal scale, similar
questions appear in access to legal remedy, recognition of workplace
injury, treatment by institutions, and the historical valuation of whose
labor and life count.

The lecture's key question is not whether people are ``naturally''
competitive. It is how rules turn some people's claims into enforceable
rights while leaving others with the cost, the debt, or the burden of
proof.

\begin{center}\rule{0.5\linewidth}{0.5pt}\end{center}

\subsection{10. Debt, Wage Slavery, and
Repair}\label{debt-wage-slavery-and-repair}

``Wage slavery'' is a polemical phrase used to criticize workers'
dependence on wages under conditions where exit is difficult. It can
illuminate coercion, low bargaining power, and the concentration of
ownership. It must not be used as if wage labor and chattel slavery were
identical: chattel slavery legally treated people as property and was
enforced through racialized violence, family separation, and the denial
of personhood. Debt bondage and peonage are also distinct systems,
though they can intersect with wage dependence and coercive labor.

Debt is both a contract and a relation of power. Who can borrow cheaply,
whose debt is forgiven, what collateral can be seized, and which
obligations survive crisis are political questions. When essential needs
require high-cost borrowing, or when a worker cannot leave without
losing housing, health care, or immigration security, ``choice'' needs
to be examined in context.

\subsubsection{Poetic justice, accountability, and
reparations}\label{poetic-justice-accountability-and-reparations}

``Poetic justice'' often means an outcome that symbolically fits an
injustice: a wrongdoer receives a fitting consequence. Symbolic
recognition can matter, but symmetry or revenge is not the same as
repair. A reparative framework asks:

\begin{enumerate}
\def\labelenumi{\arabic{enumi}.}
\tightlist
\item
  What happened, and who had the power to cause or prevent it?
\item
  What losses and continuing effects can be identified?
\item
  What forms of acknowledgment, restitution, compensation, or
  institutional change are due?
\item
  Who participates in setting the remedy and evaluating whether it
  works?
\end{enumerate}

Reparations are not simply charity. They address wrongful taking,
exclusion, or state-backed harm through material and institutional
measures. Their design may include direct payments, return of land or
property, debt relief, public investment, access to services,
truth-telling, memorialization, and guarantees against recurrence. These
mechanisms answer different harms; a single payment or apology may be
inadequate.

\subsubsection{Reconstruction: the United States and
beyond}\label{reconstruction-the-united-states-and-beyond}

U.S. Reconstruction after the Civil War (roughly 1865--1877) brought
constitutional amendments abolishing slavery, establishing birthright
citizenship, and prohibiting denial of voting rights on the basis of
race, color, or previous condition of servitude; women remained excluded
from voting. It also saw Black political participation and
institution-building. Its possibilities were sharply constrained and
then attacked through white supremacist violence, political retreat, and
the withdrawal of federal protection. Reconstruction should not be
reduced to a ``failed experiment'': its gains, limits, and reversal
reveal how rights require durable institutions and enforcement.

``Reconstruction'' also names different projects elsewhere. Postwar
rebuilding in Europe and Japan involved state capacity, investment, and
security arrangements; South Africa's post-apartheid transition paired
political change with truth-seeking and a Reconstruction and Development
Programme, while debate over land, wealth, and material redress
continued. These histories are not interchangeable. In each, ask whose
security was restored, who controlled resources, which harms were
acknowledged, and whether the project transformed the distribution of
power or mainly stabilized it.

\subsubsection{Universal basic income: promise and
peril}\label{universal-basic-income-promise-and-peril}

A universal basic income (UBI) generally means regular cash paid
individually, without a means test or work requirement. It promises
income security, administrative simplicity, greater bargaining power,
and time for care, study, or creative work. It may reduce the stigma and
cliffs that can accompany targeted benefits.

But ``UBI'' covers many different proposals. A modest payment cannot
replace health care, disability support, housing, child care, or public
services. If it substitutes for those programs, it can leave people
worse off. A poorly funded transfer could also be offset by higher
rents, prices, or cuts elsewhere. A universal payment's fiscal cost, tax
design, inflationary conditions, and distributional effects matter.
Evidence from cash-transfer trials can inform these debates, but short
pilots are not equivalent to a permanent national UBI.

Evaluate any proposal by asking: How much, how often, and to whom? What
happens to existing benefits and public services? How is it financed?
Does it improve workers' ability to refuse dangerous or degrading work?
Does it reach people without stable addresses or bank access? Does it
complement reparations and targeted support, or get offered as a
substitute for them?

\begin{center}\rule{0.5\linewidth}{0.5pt}\end{center}

\subsection{11. Synthesis: From Hoarding to Shared
Security}\label{synthesis-from-hoarding-to-shared-security}

The word ``greed'' can name a judgment, but it cannot do the whole
analysis. An adequate account connects motive, opportunity,
institutional permission, distribution of harm, and possible
alternatives:

\[
\text{harmful accumulation}
= \text{concentrated control}
+ \text{externalized costs}
+ \text{weak accountability}
\]

This is a heuristic, not a diagnostic formula. Its purpose is to direct
attention toward the social arrangements that make harm profitable,
invisible, or difficult to contest. The repair question is not simply
``How can we make individuals less greedy?'' It is ``How can we build
conditions in which safety, resources, credibility, and the power to
refuse are not hoarded?''

\begin{center}\rule{0.5\linewidth}{0.5pt}\end{center}

\subsection{12. Seminar Discussion}\label{seminar-discussion}

\begin{enumerate}
\def\labelenumi{\arabic{enumi}.}
\tightlist
\item
  In ``Known Unknowns,'' where does uncertainty come from: memory,
  Justin's account, family silence, institutional barriers, or some
  combination? Who is asked to bear the cost of not knowing?
\item
  What makes Vienna's listening ``aggressive'' in the Bureau for
  Listening sense? How would the ethical analysis change if the speaker
  had not agreed to meet or if there were no safety team?
\item
  If a synthetic voice sounds like a loved one, what would count as
  listening to that person rather than reacting to an imitation? What
  evidence and consent are missing?
\item
  How can creators and platforms distinguish permission to use a voice
  sample from permission to generate new statements in that voice?
\item
  How can we discuss toxic masculinity without treating masculinity
  itself as toxic or excusing an individual's actions as products of
  culture?
\item
  What is gained and lost when grief is described as a sequence of
  stages? How might legal accountability and grief coexist?
\item
  Humans are animals; what continuities in resource defense and
  cooperation are useful to study? Which differences arise from
  institutions, and which are exaggerated by human exceptionalism?
\item
  When is a contest actually zero-sum, and when can cooperation grow the
  pie? How would a larger pie still leave questions of ownership,
  distribution, and externalized harm?
\item
  How could epidemiologic methods identify preventable social harm
  without pathologizing a population?
\item
  Which behavioral intervention changes the environment, and which
  merely asks people at risk to change their behavior?
\item
  Can UBI increase freedom while leaving unequal ownership intact? How
  should it relate to reparations, debt relief, and universal public
  services?
\end{enumerate}

\begin{center}\rule{0.5\linewidth}{0.5pt}\end{center}

\subsection{13. Laboratory Exercise: A Harm-and-Repair
Map}\label{laboratory-exercise-a-harm-and-repair-map}

Choose a case study from a public report, a historical source, a
synthetic-media incident, or a fictional scenario. Do not use a
classmate's personal experience.

\begin{enumerate}
\def\labelenumi{\arabic{enumi}.}
\tightlist
\item
  \textbf{Define the outcome:} Specify the harm without using a vague
  label such as ``toxicity.''
\item
  \textbf{Map exposure and protection:} Identify who is more exposed,
  who has protective resources, and what evidence supports the map.
\item
  \textbf{Trace the mechanism:} Distinguish individual action, social
  norm, institutional rule, and historical condition.
\item
  \textbf{Estimate costs:} List physical, cognitive/emotional, social,
  economic, and ecological effects. Mark which cannot responsibly be
  monetized.
\item
  \textbf{Design a prevention step:} Use a behavioral or institutional
  change. State who controls it, who can opt out, and how success will
  be measured.
\item
  \textbf{Design a repair step:} Specify acknowledgment, restitution,
  compensation, service, or institutional reform. Explain who helps
  decide and how the remedy avoids replacing public obligations with
  charity.
\item
  \textbf{Stress-test the proposal:} Ask who might be excluded,
  surveilled, burdened, or harmed by the intervention.
\item
  \textbf{For a voice-cloning case:} Separate the audio's apparent
  identity from verified source, consent, context, and intended
  audience. Propose a verification path that does not require trusting
  the disputed recording.
\end{enumerate}

\textbf{Deliverable:} A one-page causal map and a short explanation of
one uncertainty that remains. The goal is not to produce a perfect
measure of suffering; it is to make the reasoning, evidence, and limits
visible.

\begin{center}\rule{0.5\linewidth}{0.5pt}\end{center}

\subsection{References}\label{references}

\begin{enumerate}
\def\labelenumi{\arabic{enumi}.}
\tightlist
\item
  American Psychological Association. (2018). \emph{Guidelines for
  Psychological Practice with Boys and Men}.
  \url{https://www.apa.org/about/policy/boys-men-practice-guidelines.pdf}
\item
  Bureau for Listening. ``Aggressive Listening.''
  \url{https://bureauforlistening.com/agressive-listening/}
\item
  Connell, R. W. (1995). \emph{Masculinities}. University of California
  Press.
\item
  DeKornfeld, A. (host). (2022). ``902: The Missing Minutes.''
  \emph{This American Life}.
  \url{https://www.thisamericanlife.org/902/transcript}
\item
  Du Bois, W. E. B. (1935). \emph{Black Reconstruction in America,
  1860--1880}. Harcourt, Brace and Company.
\item
  Foner, Eric. (1988). \emph{Reconstruction: America's Unfinished
  Revolution, 1863--1877}. Harper \& Row.
\item
  Friston, Karl. (2010). ``The free-energy principle: a unified brain
  theory?'' \emph{Nature Reviews Neuroscience}, 11(2), 127--138.
\item
  Kahneman, Daniel, and Amos Tversky. (1979). ``Prospect Theory: An
  Analysis of Decision under Risk.'' \emph{Econometrica}, 47(2),
  263--291.
\item
  Kupers, Terry A. (2005). ``Toxic masculinity as a barrier to mental
  health treatment in prison.'' \emph{Journal of Clinical Psychology},
  61(6), 713--724.
\item
  Klass, Dennis, Phyllis R. Silverman, and Steven L. Nickman,
  eds.~(1996). \emph{Continuing Bonds: New Understandings of Grief}.
  Taylor \& Francis.
\item
  Rothman, Kenneth J., Sander Greenland, and Timothy L. Lash. (2008).
  \emph{Modern Epidemiology}, 3rd ed.~Lippincott Williams \& Wilkins.
\item
  Thaler, Richard H., and Cass R. Sunstein. (2008). \emph{Nudge:
  Improving Decisions About Health, Wealth, and Happiness}. Yale
  University Press.
\item
  United Nations Department of Economic and Social Affairs. ``Universal
  Basic Income Policies and Their Potential for Poverty Reduction.''
  \url{https://www.un.org/development/desa/dspd/2020/02/13609/}
\item
  World Health Organization. ``Social determinants of health.''
  \url{https://www.who.int/health-topics/social-determinants-of-health}
\item
  American Psychiatric Association. (2022). \emph{Diagnostic and
  Statistical Manual of Mental Disorders}, 5th ed., text rev.
  \url{https://doi.org/10.1176/appi.books.9780890425787}
\item
  Blackmon, Douglas A. (2008). \emph{Slavery by Another Name: The
  Re-Enslavement of Black Americans from the Civil War to World War II}.
  Doubleday.
\item
  Darity, William A., Jr., and A. Kirsten Mullen. (2020). \emph{From
  Here to Equality: Reparations for Black Americans in the Twenty-First
  Century}. University of North Carolina Press.
\item
  Hoynes, Hilary, and Jesse Rothstein. (2019). ``Universal Basic Income
  in the United States and Advanced Countries.'' \emph{Annual Review of
  Economics}, 11, 929--958.
  \url{https://doi.org/10.1146/annurev-economics-080218-030237}
\item
  Parker, Geoffrey A. (1974). ``Assessment strategy and the evolution of
  fighting behaviour.'' \emph{Journal of Theoretical Biology}, 47(1),
  223--243. \url{https://doi.org/10.1016/0022-5193(74)90111-8}
\item
  \emph{Unlearning Insecurity: Active Inference, Bayesian Precision
  Weighting, and Cognitive Restructuring.} Companion notebook, Chaos
  Hearing.
  \href{../../../hpc/chaoshearing/notebooks/unlearning_insecurity_active_inference.ipynb}{Local
  file}.
\item
  Bliss, Shepherd. (1995). ``Mythopoetic Men's Movement.'' In Michael S.
  Kimmel (Ed.), \emph{The Politics of Manhood: Profeminist Men Respond
  to the Mythopoetic Men's Movement (and the Mythopoetic Leaders
  Answer)}. Temple University Press.
\item
  Dower, John W. (1999). \emph{Embracing Defeat: Japan in the Wake of
  World War II}. W. W. Norton.
\item
  Judt, Tony. (2005). \emph{Postwar: A History of Europe Since 1945}.
  Penguin Press.
\item
  Truth and Reconciliation Commission of South Africa. (1998).
  \emph{Truth and Reconciliation Commission of South Africa Report}.
\item
  Government of South Africa. (1994). \emph{White Paper on
  Reconstruction and Development}.
\item
  Worden, J. William. (2018). \emph{Grief Counseling and Grief Therapy:
  A Handbook for the Mental Health Practitioner}, 5th ed.~Springer
  Publishing Company.
\item
  Federal Trade Commission. (2023). ``Scammers use AI to enhance their
  family emergency schemes.''
  \url{https://consumer.ftc.gov/consumer-alerts/2023/03/scammers-use-ai-enhance-their-family-emergency-schemes}
\item
  National Institute of Standards and Technology. (2024). \emph{Reducing
  Risks Posed by Synthetic Content: An Overview of Technical Approaches
  to Digital Content Transparency}. NIST AI 100-4.
  \url{https://nvlpubs.nist.gov/nistpubs/ai/NIST.AI.100-4.pdf}
\item
  Nowak, Martin A. (2006). ``Five rules for the evolution of
  cooperation.'' \emph{Science}, 314(5805), 1560--1563.
  \url{https://doi.org/10.1126/science.1133755}
\end{enumerate}

\end{document}
